\documentclass[12pt]{article}

% Match weeks 2--4 using fonts available in the local TeX installation.
\usepackage[T1]{fontenc}
\usepackage{times}
\usepackage{microtype}
\usepackage{amsmath}
\usepackage[parfill]{parskip}
\setlength{\parindent}{0pt}
\setlength{\parskip}{2ex plus 0.2ex minus 0.2ex}
\usepackage[letterpaper,margin=1in]{geometry}
\usepackage[hidelinks]{hyperref}

\begin{document}

\begin{flushleft}
\textbf{Week 5: Topic Models and Mixed-Membership Models} \\
Jaisel Singh (js6897) \\
October 14, 2026
\end{flushleft}

\vspace{0.1in}

What stood out to me in Blei's review was the shift from clustering a document to
clustering its words. A topic is a distribution over a fixed vocabulary, every
document in the collection draws on the same topics, and what distinguishes one
document from another is the proportion in which it uses them. Because a topic is
chosen separately for each word, a single document can express several topics at
once. That is what mixed membership means, and it is what a model assigning one
cluster per document cannot represent. The generative story is short: draw topic
proportions for a document, then for each word draw a topic assignment and draw a
word from the corresponding topic. I noticed the Dirichlet-categorical structure
sitting underneath both of those draws.

What I found most useful was how cleanly the model separates from the computation.
The words are observed and the topic structure (the topics, the per-document
proportions, and the per-word assignments) is hidden, so fitting the model means
computing a posterior over that hidden structure. Blei describes this as reversing
the generative process. The obstacle is the evidence in the denominator, which sums
the joint distribution over every way of assigning each observed word to a topic,
and that count grows exponentially with the collection. This is the same wall we
hit last week, and the responses are the same two families: sampling, usually Gibbs
sampling, or variational approximation. It makes the MCMC material feel less like a
separate topic and more like the machinery these models need in order to be fit at all.

Boyd-Graber, Hu, and Mimno's survey complicated my sense of what a good fit means
here. They report that held-out likelihood can correlate negatively with how
interpretable readers find a model, since a model with more topics can predict
held-out words better while being harder for a person to read. Automated coherence
scores stand in for human judgment, but they do not establish that a model suits a
particular application. My remaining question is how to set the number of topics and
the Dirichlet hyperparameters when the quantity I can compute is not the quantity I
actually care about.

\vspace{0.05in}
{\footnotesize
\textit{Readings:} Blei (2012), \href{https://www.cs.columbia.edu/~blei/fogm/readings/Blei2012.pdf}{\emph{Probabilistic Topic Models}};
Boyd-Graber et al.\ (2017), \href{https://www.nowpublishers.com/article/Details/INR-030}{\emph{Applications of Topic Models}}.
\par}

\end{document}

%%% Local Variables:
%%% mode: latex
%%% TeX-master: t
%%% End:
